%Version 3.1 December 2024 -- Springer Nature sn-jnl template
%%%%%%%%%%%%%%%%%%%%%%%%%%%%%%%%%%%%%%%%%%%%%%%%%%%%%%%%%%%%%%%%%%%%%%
\documentclass[pdflatex,sn-mathphys-num]{sn-jnl}% Math and Physical Sciences Numbered Reference Style

\usepackage{graphicx}%
\usepackage{multirow}%
\usepackage{amsmath,amssymb,amsfonts}%
\usepackage{amsthm}%
\usepackage{mathrsfs}%
\usepackage[title]{appendix}%
\usepackage{xcolor}%
\usepackage{textcomp}%
\usepackage{manyfoot}%
\usepackage{booktabs}%
\usepackage{algorithm}%
\usepackage{algorithmicx}%
\usepackage{algpseudocode}%
\usepackage{listings}%
\usepackage{placeins}%

\theoremstyle{thmstyleone}%
\newtheorem{theorem}{Theorem}%
\newtheorem{proposition}[theorem]{Proposition}%

\theoremstyle{thmstyletwo}%
\newtheorem{example}{Example}%
\newtheorem{remark}{Remark}%

\theoremstyle{thmstylethree}%
\newtheorem{definition}{Definition}%

\raggedbottom

\makeatletter
\@ifpackageloaded{hyperref}{%
  \hypersetup{
    pdftitle={Does Protein-Language-Model Scale Improve Zero-Shot Identification of Kinase-Inhibitor Resistance?},
    pdfauthor={Mohamed Samiulla Khan},
    pdfsubject={Zero-shot protein language model scoring of kinase-inhibitor resistance},
    pdfkeywords={protein language models, kinase inhibitor resistance, ESM-2, deep mutational scanning}
  }
}{}
\makeatother

\begin{document}

\title[PLM scale and kinase-inhibitor resistance]{Does Protein-Language-Model Scale Improve Zero-Shot Identification of Kinase-Inhibitor Resistance? A Multi-Kinase Test Against Non-Learned Baselines}

\author*[1]{\fnm{Mohamed Samiulla} \sur{Khan}}\email{mkhan@mec.edu.om}

\affil[1]{\orgdiv{Department of Computing and Electronics Engineering}, \orgname{Middle East College}, \orgaddress{\city{Muscat}, \country{Oman}}}

\abstract{Protein language models (PLMs) trained on evolutionary sequence data can predict mutation effects without task-specific training, and larger PLMs generally perform better on standard fitness-prediction benchmarks. Whether this scale advantage extends to a pharmacologically specific quantity -- an inhibitor-resistance mutation's effect on a particular drug, rather than general protein fitness -- has not been established. We test zero-shot masked-marginal scoring from ESM-1b and three ESM-2 sizes (150M, 650M, 3B parameters) against measured resistance data across three tyrosine kinases (MET, EGFR, ABL1), benchmarked against non-learned baselines including two structure-derived distances from real inhibitor-bound crystal structures. For MET (5{,}434 quality-controlled variants from an original panel of 5{,}764, evaluated against nine inhibitors), ESM-1b significantly outperforms every ESM-2 size and all eight non-learned baselines, with 95\% bootstrap confidence intervals excluding zero for every comparison; scale improves performance monotonically within the ESM-2 family alone without closing the gap to ESM-1b. These claims are scoped to the nine inhibitors tested. All four model sizes also correlate substantially better with untreated, drug-agnostic fitness than with drug-treated resistance, extending a contrast Estevam et al.\ reported for ESM-1b alone to the full model family. For EGFR and ABL1, matched to ChEMBL bioactivity data, no predictor reaches significance after correction across 21 exact-permutation tests, a result that persists after panel expansion, leave-one-out analysis, and paired comparison against every baseline; a post-hoc power analysis shows the available sample sizes ($n=6$--$15$) give only 8--13\% power to detect a MET-sized effect. We interpret this null result not as evidence that scale is irrelevant to drug-specific resistance, but as a demonstration that deep mutational scanning and curated bioactivity databases offer fundamentally different statistical power for this comparison.}

\keywords{protein language models, zero-shot variant effect prediction, kinase inhibitor resistance, ESM-2, drug-target affinity, deep mutational scanning}

\maketitle

\section{Introduction}\label{sec1}

Small-molecule kinase inhibitors are among the most successful classes of targeted cancer therapy. Imatinib's original phase-1 trial in BCR-ABL1-driven chronic myeloid leukemia \citep{Druker2001} and gefitinib's IPASS trial in EGFR-mutant pulmonary adenocarcinoma \citep{Mok2009} exemplify the durable, high-response-rate benefit that made kinase inhibitors a template for targeted oncology more broadly; by 2022, 68 small-molecule kinase inhibitors targeting roughly two dozen distinct kinases had received FDA approval \citep{Roskoski2022}. That clinical benefit, however, is routinely eroded by acquired resistance mutations in the kinase domain itself. The gatekeeper mutation T790M and the covalent-anchor mutation C797S limit the durability of EGFR-directed therapy in non-small-cell lung cancer \citep{Kobayashi2005,Thress2015}; the gatekeeper mutation T315I remains the most clinically consequential resistance mutation in BCR-ABL1-driven chronic myeloid leukemia, first identified as a cause of clinical STI-571 (imatinib) resistance and later shown to require a dedicated pan-BCR-ABL inhibitor (ponatinib) to overcome \citep{Gorre2001,OHare2009}. Anticipating which mutations will confer resistance to a given inhibitor -- ideally before they are observed in the clinic -- is a long-standing goal that current sequence-based computational tools do not reliably achieve.

Protein language models (PLMs) trained by masked-language-modeling on large evolutionary sequence corpora have shown that a model never trained on any functional or fitness label can still predict, in zero-shot fashion, how a mutation is likely to affect protein function \citep{Meier2021,Rives2021}. This zero-shot capability is validated primarily against deep mutational scanning (DMS) datasets and other general fitness-proxy benchmarks, where it captures broad evolutionary and structural constraint reasonably well. Whether this same zero-shot signal is informative about a substantially narrower and more clinically specific quantity -- how a mutation shifts a protein's affinity for one particular small-molecule inhibitor -- is a separate empirical question. A model with no access to ligand information at inference time has no obvious mechanism for capturing drug-specific resistance except insofar as drug-specific effects happen to correlate with general destabilization or loss of evolutionary constraint, and there is no a priori reason to assume they do.

The most direct evidence bearing on this question comes from Estevam et al.\ \citep{Estevam2025}, who performed a deep mutational scan of 5{,}764 MET kinase-domain variants against a panel of nine clinically relevant inhibitors and reported that zero-shot ESM-1b scores correlate only weakly with drug-treated resistance fitness (pooled $r \approx 0.28$), markedly lower than the correlation with untreated, drug-agnostic fitness ($r \approx 0.50$) computed on the same variants. Their proposed explanation is that drug-specific resistance is not an evolutionary pressure the model has seen during pretraining, so the zero-shot signal -- while real -- is diluted when applied to a drug-specific readout. This is a single-kinase, single-model-generation result: it leaves open whether the weak correlation is a property of MET specifically, of ESM-1b specifically, or of the zero-shot approach in general regardless of model scale or target kinase.

Model scale is a natural candidate explanation to rule in or out. Successive generations of ESM models substantially improve on general protein-fitness and structure-prediction benchmarks with increasing parameter count \citep{Lin2023}, and it would be a reasonable prior that a larger, more capable model narrows the gap Estevam et al.\ observed. It is an equally reasonable prior that it would not: if the underlying limitation is that drug-specific resistance is simply not encoded in evolutionary sequence statistics at all, then no amount of additional scale within the same training paradigm should be expected to help, and a well-powered structural or substitution-matrix baseline computed with no learned model at all should perform just as well. Distinguishing between these two possibilities -- and doing so across more than one kinase, since a single-kinase result cannot establish whether either pattern generalizes -- is the aim of this study.

We test zero-shot masked-marginal scoring \citep{Meier2021} from ESM-1b and three sizes of ESM-2 (150M, 650M, and 3B parameters) against real, measured resistance data for three kinases spanning two structural classes -- MET and EGFR (receptor tyrosine kinases) and ABL1 (a non-receptor tyrosine kinase) -- and benchmark every model against non-learned baselines: a BLOSUM62 substitution score \citep{Henikoff1992}, and two structure-derived distances computed from real inhibitor-bound crystal structures rather than approximated. For MET, where the deep-mutational-scanning design of Estevam et al.\ \citep{Estevam2025} provides thousands of variants against nine inhibitors, this comparison is well powered and yields a clear answer: scale does not close the gap between ESM-1b and ESM-2 for drug-specific resistance prediction, at any tested size. For EGFR and ABL1, the same comparison is not statistically resolvable at the sample sizes available even after literature-directed panel expansion, a limitation we quantify directly with permutation tests, leave-one-out analysis, and a power calculation (\S\ref{subsec2-3}). We report both results, and the sharp difference in what each is able to establish, as evidence for a broader methodological point developed in \S\ref{sec4}: deep-mutational-scanning-derived and bioactivity-database-derived resistance datasets offer fundamentally different statistical power for this class of question.

\section{Related Work}\label{secRW}

This study sits at the intersection of three literatures that have so far remained largely separate: supervised drug-target affinity prediction, structure-based variant-effect prediction, and unsupervised zero-shot variant-effect prediction. None of the three, on its own, has directly tested what we test here -- whether zero-shot, sequence-only PLM scoring, with no drug-specific training signal of any kind, tracks measured resistance to a \emph{specific} inhibitor, across more than one kinase and more than one model scale.

\textbf{Supervised drug-target affinity (DTA) prediction.} The Davis kinase-inhibitor selectivity panel \citep{Davis2011} and the KIBA dataset \citep{Tang2014}, which integrates heterogeneous $K_i$/$K_d$/IC$_{50}$ measurements across a large kinase-compound matrix, established the standard large-scale benchmarks for kinase-inhibitor binding-affinity prediction. DeepDTA \citep{Ozturk2018}, GraphDTA \citep{Nguyen2021}, and the DeepPurpose toolkit \citep{Huang2020} are representative supervised architectures built on these benchmarks: convolutional, graph-neural-network, and modular encoder-based models, respectively, that learn a compound-target affinity function end-to-end from thousands of labeled binding measurements. These models are trained directly on affinity labels and answer a fundamentally different question from ours -- given substantial supervision, how accurately can affinity be predicted -- whereas we ask whether a model with \emph{zero} exposure to any affinity label, drug-specific or otherwise, retains any signal about drug-specific resistance at all. A supervised DTA model that performs well on DAVIS or KIBA provides no evidence either way on this question, since it has seen the very labels we are asking whether a zero-shot model can do without.

\textbf{Structure prediction and general variant pathogenicity.} AlphaFold's near-experimental-accuracy structure prediction \citep{Jumper2021} has enabled structure-aware downstream tools including AlphaMissense \citep{Cheng2023}, which combines predicted structure with evolutionary sequence context to classify missense variants as pathogenic or benign genome-wide. AlphaMissense and related structure-informed classifiers are calibrated against general pathogenicity (e.g.\ ClinVar labels), not against resistance to a specific small-molecule inhibitor, and so do not directly address drug-specific resistance either. We deliberately used experimentally resolved, inhibitor-bound crystal structures (\S\ref{subsec3-3}) rather than predicted structures for our own structural baselines, precisely to isolate whether real, ground-truth information about the ligand-binding pocket carries drug-specific resistance signal, independent of any uncertainty a structure-prediction step would introduce.

\textbf{Zero-shot and generative variant-effect prediction.} The general principle this study builds on -- that models trained only on evolutionary sequence statistics, with no fitness or functional label, can predict mutation effects -- predates ESM. Alignment-based generative models including DeepSequence \citep{Riesselman2018} and EVmutation \citep{Hopf2017}, and their more recent extension EVE \citep{Frazer2021}, established this using variational and pairwise-coupling models over multiple sequence alignments rather than large transformer architectures. Tranception \citep{Notin2022} extended zero-shot scoring to autoregressive transformers with retrieval augmentation, evaluated on the ProteinGym benchmark suite; TAPE \citep{Rao2019} established a broader benchmark suite for learned protein representations, including some structure- and stability-related tasks but no drug-resistance task. All of this work validates zero-shot and generative variant-effect models against general fitness, stability, or pathogenicity endpoints. None evaluates whether the same zero-shot signal is informative about resistance to one specific small-molecule inhibitor, which is the endpoint Estevam et al.\ \citep{Estevam2025} first showed is only weakly captured by ESM-1b for MET, and which is the question this study extends across model scale and kinase identity.

\textbf{Kinase-resistance-adjacent PLM applications.} Closely related but methodologically distinct work exists specifically in the kinase-resistance and PLM-for-binding space. Pan et al.\ \citep{Pan2026} benchmark structure/docking-based scoring tools for mutation-induced protein-ligand affinity changes in the absence of experimental structures; this is a structure-based scoring-function comparison, not a zero-shot, sequence-only PLM evaluation, and is not kinase-resistance-specific. Gurusinghe et al.\ \citep{Gurusinghe2025} (ProBASS) use ESM-2 and ESM-IF1 embeddings to predict $\Delta\Delta G_{\text{bind}}$, but for protein--protein interactions and via supervised fine-tuning rather than zero-shot scoring of drug binding. Singh et al.\ \citep{Singh2023} (ConPLex) use PLM embeddings for general drug--target interaction prediction, experimentally validated on real kinase--drug pairs, but do not address mutation-specific resistance sensitivity. Wu et al.\ \citep{Wu2025} curate a modification-aware extension of the DAVIS benchmark spanning ten kinases including EGFR and ABL1, but evaluate supervised affinity-model generalization to unseen mutations rather than an off-the-shelf PLM's intrinsic zero-shot signal. None of these studies test zero-shot, sequence-only PLM scoring against real measured kinase-inhibitor resistance across multiple kinases and multiple model scales, which is the gap this study addresses.

\section{Results}\label{sec2}

\subsection{Partial reproduction of the MET zero-shot scoring pipeline}\label{subsec2-1}

Before extending the analysis to new kinases, we compared our own zero-shot scoring implementation against Estevam et al.'s published per-mutation ESM-1b scores for all 5{,}434 quality-controlled MET kinase-domain variants. Our masked-marginal reimplementation correlated closely with their published scores (Spearman $\rho = 0.9199$), and this figure was unchanged across a second, independent execution of the identical pipeline. We are explicit about what this second run does and does not establish: it confirms the pipeline is deterministic and repeatable, not that it is correct in the sense of matching Estevam et al.'s scores exactly, and $\rho = 0.9199$ remains below our a priori target of $\rho \geq 0.95$. We therefore describe this as a partial reproduction rather than a full replication. Candidate sources of the residual gap that we have not exhaustively ruled out include the exact construct sequence used at inference time, residue-numbering conventions, the masked- versus wildtype-marginal scoring convention, the precise model checkpoint and library version, floating-point precision, sequence-preprocessing choices, and mutation-filtering criteria. We tested one of these candidates directly -- the masked- versus wildtype-marginal scoring convention (\S\ref{subsec2-4}) -- and found it explains only a small fraction of the gap; the remaining candidates are untested and are listed explicitly in \S\ref{subsec4-1} as open items rather than assumed to be immaterial.

\subsection{Untreated fitness correlations exceed treated correlations across all four models}\label{subsec2-1b}

The central mechanistic argument motivating this study -- that a PLM pretrained only on evolutionary sequence data should track general, drug-agnostic protein fitness far better than drug-specific resistance, because only the former resembles a pressure the model has seen during pretraining -- rests on comparing each model's correlation with untreated (DMSO) fitness against its correlation with drug-treated fitness. Estevam et al.\ \citep{Estevam2025} report this contrast for their own ESM-1b scores: untreated $r \approx 0.50$ versus drug-treated $r \approx 0.28$. We recomputed this contrast for our own from-scratch scorers across all four model sizes on the untreated (DMSO) condition ($n=4{,}595$ variants with complete untreated scores): ESM-1b ($\rho=0.500$), ESM-2-150M ($\rho=0.469$), ESM-2-650M ($\rho=0.500$), and ESM-2-3B ($\rho=0.500$). Every model showed a substantially higher correlation with untreated fitness than with drug-treated fitness (Fisher-z-averaged drug-treated correlations from \S\ref{subsec2-2}; gaps of 0.246--0.270 across the four models). We tested this gap against a null of zero with a bootstrap that resamples mutations once per iteration and applies the same resampled indices to both the untreated and treated correlations (2{,}000 iterations, $n=4{,}595$ shared mutations; full procedure in \S\ref{subsec3-5}): the 95\% bootstrap CI excluded zero for every model (ESM-1b: $0.242$, 95\% CI $[0.216, 0.268]$; ESM-2-150M: $0.255$, CI $[0.227, 0.282]$; ESM-2-650M: $0.267$, CI $[0.241, 0.294]$; ESM-2-3B: $0.260$, CI $[0.234, 0.287]$; all $p \leq 0.001$, the resolution limit of a 2{,}000-iteration bootstrap). This supports the interpretation that the untreated-versus-treated performance gap is not specific to ESM-1b: all four model sizes track general, drug-agnostic MET fitness far better than they track drug-specific resistance, and the gap size itself does not track model scale in any obvious direction. We note that three of the four untreated correlations round to an identical $\rho=0.500$ at three decimal places; the underlying values are close but not identical, and we report the values as computed without further rounding adjustment.

\subsection{MET: ESM-1b outperforms every tested ESM-2 size and every non-learned baseline}\label{subsec2-2}

Across all nine MET inhibitors, zero-shot ESM-1b scores correlated with drug-treated resistance fitness with a pooled Spearman $\rho = 0.257$, exceeding every non-learned baseline tested: BLOSUM62 ($\rho = 0.175$), pocket volume ($\rho = 0.063$), distance to the ATP pocket ($\rho = 0.054$), hydrophobicity ($\rho = 0.028$), distance to the reference ligand ($\rho = -0.018$), a point mutation stability estimate ($\rho = -0.011$), polarity ($\rho \approx 0.000$), and an aggregate stability estimate ($\rho = -0.169$) (Figure~\ref{fig2}). This pooled comparison is explicitly a descriptive summary, not the primary inferential analysis: it repeats each of the 5{,}434 mutations nine times (once per drug), so the $n=9\times5{,}434$ observations underlying it are not independent, and a pooled correlation can in principle be shaped by between-drug scale differences or drug-specific offsets in ways a per-drug analysis is not. We therefore additionally ran a formal paired significance test of ESM-1b against each of the eight non-learned baselines using the same shared-mutation cluster bootstrap procedure as the ESM-1b-versus-ESM-2 comparison below (2{,}000 iterations, one shared mutation-index set per iteration applied to ESM-1b and every baseline, Fisher-z-averaged across the nine drugs before recording each iteration's difference; full procedure in \S\ref{subsec3-5}). ESM-1b significantly outperformed all eight baselines: every 95\% bootstrap CI on the Fisher-z-averaged difference excluded zero, and after Benjamini-Hochberg correction all adjusted bootstrap $p$-values were at or below the resolution limit of the 2{,}000-iteration bootstrap ($p \leq 0.001$): BLOSUM62 (mean difference $+0.079$, 95\% CI $[0.051, 0.107]$), distance to the ATP pocket ($+0.201$, CI $[0.174, 0.228]$), distance to the reference ligand ($+0.263$, CI $[0.240, 0.285]$), the point-mutation stability estimate ($+0.263$, CI $[0.227, 0.299]$), the aggregate stability estimate ($+0.418$, CI $[0.376, 0.461]$), pocket volume ($+0.258$, CI $[0.231, 0.283]$), hydrophobicity ($+0.253$, CI $[0.227, 0.280]$), and polarity ($+0.261$, CI $[0.235, 0.288]$).

Per-drug correlations ranged from $\rho = 0.039$ (capmatinib) to $\rho = 0.414$ (savolitinib) for ESM-1b, roughly an eleven-fold spread that is itself informative: a single pooled correlation obscures substantial per-drug heterogeneity (Figure~\ref{fig1}). At every one of the nine inhibitors, ESM-1b's correlation with drug-treated fitness was numerically higher than or comparable to all three ESM-2 sizes. A paired Wilcoxon signed-rank test across the nine drugs confirmed this is not incidental: ESM-1b significantly outperformed ESM-2-650M and ESM-2-3B (9/9 drugs, $p = 0.0039$ for both comparisons) and ESM-2-150M (8/9 drugs, $p = 0.0078$). This test treats the nine per-drug correlations as nine independent paired observations, but all nine are computed on the same underlying panel of 5{,}434 mutations, so the sampling errors of the nine correlations are not strictly independent in the way a standard paired Wilcoxon test assumes. We therefore additionally ran a shared-mutation cluster bootstrap across the nine fixed drug conditions (2{,}000 resamples; each resample draws a set of \emph{mutations} with replacement from the full 5{,}434-mutation panel, applies that same resampled index set to all nine drugs, and recomputes all nine per-drug correlations and the Fisher-z-averaged summary within each resample -- full procedural detail in \S\ref{subsec3-5}), which propagates the shared-panel dependence into the uncertainty estimate directly rather than assuming independence across drugs; because the nine drug conditions themselves are held fixed and only mutations are resampled, this quantifies uncertainty from the finite mutation panel, not from the broader population of possible drugs. This confirmed the Wilcoxon result: ESM-1b's advantage over ESM-2-150M (mean Fisher-z-averaged difference $+0.040$, 95\% CI $[0.025, 0.055]$), ESM-2-650M ($+0.023$, CI $[0.011, 0.035]$), and ESM-2-3B ($+0.015$, CI $[0.005, 0.027]$) all excluded zero. Scale improved performance monotonically \emph{within} the ESM-2 family alone (cluster-bootstrap Fisher-z-averaged $\rho$: 150M $= 0.214$, 650M $= 0.231$, 3B $= 0.239$), but this improvement never closed the gap to ESM-1b (Fisher-z-averaged $\rho = 0.254$).

\begin{figure}[!htb]
\centering
\includegraphics[width=\textwidth]{figure1_met_per_drug_model_comparison.pdf}
\caption{MET kinase domain: zero-shot score versus drug-treated resistance fitness, per inhibitor and per model. Inhibitor abbreviations (as ordered on the $x$-axis): Savo = savolitinib, Tepo = tepotinib, Cabo = cabozantinib, Glu = glumetinib, Crizo = crizotinib, NVP = the NVP compound described in \S\ref{subsec3-1}, Mere = merestinib, Gle = glesatinib, Camp = capmatinib. ESM-1b (dark blue) numerically exceeds all three ESM-2 sizes at every inhibitor; a paired Wilcoxon signed-rank test confirms this is significant against ESM-2-650M and ESM-2-3B (9/9 drugs, $p=0.0039$) and against ESM-2-150M (8/9 drugs, $p=0.0078$). Source data: \texttt{met\_full\_model\_drug\_correlation\_table.csv}.}\label{fig1}
\end{figure}

\begin{figure}[!htb]
\centering
\includegraphics[width=0.75\textwidth]{figure2_met_baseline_comparison_v2.pdf}
\caption{MET: pooled Spearman correlation of zero-shot ESM-1b and eight non-learned baselines with drug-treated resistance fitness ($n=9$ drugs $\times$ 5{,}434 variants, each mutation repeated once per drug and therefore not independently sampled). ESM-1b is numerically higher than every non-learned baseline tested, and a shared-mutation cluster bootstrap confirms the advantage against all eight baselines: every 95\% CI excludes zero (Benjamini-Hochberg-adjusted $p \leq 0.001$ for every comparison; \S\ref{subsec2-2}, \S\ref{subsec3-5}). Read alongside, not instead of, the per-drug analysis in Figure~\ref{fig1}. Source data: \texttt{met\_baseline\_comparison.csv}, \texttt{met\_paired\_esm1b\_vs\_baselines\_v2.csv}.}\label{fig2}
\end{figure}

\FloatBarrier

\subsection{EGFR and ABL1: no predictor -- learned or non-learned -- reaches significance, and the panels are severely underpowered}\label{subsec2-3}

To test whether the MET finding generalizes, we assembled a clinically characterized resistance-mutation panel for EGFR (a cascade of L858R, L858R+T790M, and L858R+T790M+C797S, reflecting that compound mutations rather than isolated single mutations dominate real EGFR resistance data, plus the independent activating mutations L861Q, G719C, and G719S) and ABL1 (13 among the most frequently documented BCR-ABL1 kinase-domain resistance mutations in the clinical literature), matched against ChEMBL bioactivity data.

A methodological refinement was necessary here. An initial pairing of wild-type and mutant activity records used compound identity and assay type alone; because a single compound-target pair can show substantially different measured potency across different publications and assay formats even for nominally identical assay types (in our data, one compound-target pair varied by up to 190-fold across three assay formats within a single publication), this construction conflated assay heterogeneity with the true mutation effect. We rebuilt the pairing to require agreement on compound, assay type, source publication, and assay format, aggregating first within compound and then across compounds. We refer to the resulting comparisons throughout as \emph{document- and format-matched observational pairs} rather than as matched experiments: they still do not guarantee an identical assay protocol, cell or biochemical system, protein construct, incubation time, or laboratory conditions between the wild-type and mutant measurements, only agreement on compound, assay type, publication, and reported assay format. This stricter construction changed several mutations' estimated effect sizes substantially (for example, ABL1 T315I's estimated effect roughly halved and EGFR L858R's roughly doubled relative to the looser construction), while leaving the reliable mutation count essentially unchanged (9 of 13 ABL1 mutations and 6 of 6 EGFR mutations/steps met a minimum-matched-pairs reliability threshold, for 15 of 19 total).

On this corrected dataset, no predictor reached statistical significance for ABL1 ($n=9$) or EGFR ($n=6$) -- the two within-kinase groupings we treat as primary -- across all four PLM sizes and all three non-learned baselines, using exact permutation tests (362{,}880 permutations for ABL1, 720 for EGFR) rather than the asymptotic Spearman $t$-approximation, which is unreliable at these sample sizes (14 tests total; Figure~\ref{fig3}, Figure~\ref{figScatter}). We also report a pooled \emph{Combined ($n=15$)} grouping for comparability with our original analysis plan, but treat it as a secondary sensitivity check rather than a primary inference: ABL1 and EGFR predictors sometimes carry opposite signs for the same predictor (e.g.\ ESM-1b: ABL1 $\rho=+0.233$ versus EGFR $\rho=-0.543$), so pooling assumes the two kinases' mutation-score and $\Delta\mathrm{p}X$ distributions are directly rank-comparable, an assumption kinase identity could easily violate. One nominally significant result before correction (EGFR, ESM-2-150M, $\rho = -0.943$) had an asymptotic $p=0.005$ but an exact permutation $p=0.017$ (12 of 720 permutations at least as extreme); this did not survive Benjamini-Hochberg correction across the full 21-test battery under either method (adjusted $p=0.105$ asymptotic, $p=0.350$ permutation-based) \citep{BenjaminiHochberg1995}. In the pooled Combined grouping, ESM-2-3B's raw permutation $p=0.053$ likewise did not survive correction (adjusted $p=0.398$). No test in the 21-test battery survived correction under either method.

We additionally ran three robustness checks specifically aimed at this sample-size concern. First, we compared ESM-1b directly against every other predictor -- including the three non-learned baselines individually, not only the other PLM sizes -- via paired bootstrap on the difference in Spearman $\rho$ (20{,}000 resamples per comparison); every one of the 18 paired comparisons (6 alternative predictors $\times$ 3 groupings) had a 95\% CI on the difference that included zero, so ESM-1b's numerically higher pooled MET performance does not carry over to a statistically confirmed pairwise advantage on EGFR/ABL1 for any predictor. Second, a leave-one-out (jackknife) analysis showed that removing a single mutation can swing $\rho$ by as much as 0.72 (EGFR, distance to ATP pocket) and can flip its sign entirely in several cases (e.g.\ ABL1 ESM-1b: $\rho$ ranges from $-0.10$ to $+0.50$ across the nine leave-one-out subsets); Figure~\ref{figScatter} shows the underlying per-mutation data with mutation labels so this instability is visible directly rather than only summarized. Third, a post-hoc power analysis (simulated bivariate-normal data, Spearman test, $\alpha=0.05$, 20{,}000 replicates per cell) shows that at $n=6$ (EGFR) or $n=9$ (ABL1), statistical power to detect an effect the size of the pooled MET ESM-1b correlation ($\rho\approx0.26$) is only 8--10\%, rising to just 13\% even at the pooled $n=15$; reaching conventional 80\% power at that effect size would require approximately 117 mutations under a simple single-correlation design. These three checks converge on the same conclusion: the EGFR/ABL1 null result is not evidence that no such effect exists, but a direct, quantified demonstration that the available panels cannot resolve the question either way.

Two patterns are worth flagging as hypothesis-generating only, given none were statistically distinguishable from chance at this sample size: unlike MET, the one nominally significant EGFR signal (before correction) favored ESM-2 over ESM-1b, the opposite ranking from the MET result; and the ABL1 structural baseline (distance to the ATP pocket, $\rho = 0.552$) numerically exceeded every PLM tested.

\begin{figure}[!htb]
\centering
\includegraphics[width=0.85\textwidth]{figure3_egfr_abl1_forest_plot_v4.pdf}
\caption{EGFR/ABL1: Spearman correlation point estimates with 95\% bootstrap confidence intervals, for all seven predictors (four PLM sizes, BLOSUM62, distance to reference ligand, distance to ATP pocket). ABL1 ($n=9$) and EGFR ($n=6$), shaded bands, are the primary within-kinase analyses; Combined ($n=15$) is a secondary pooled sensitivity check only (\S\ref{subsec2-3}). Open circles mark the two of 21 confidence intervals that exclude zero (EGFR ESM-2-150M; combined ESM-2-3B); neither survives Benjamini-Hochberg correction for the full battery of 21 exact-permutation tests. Source data: \texttt{egfr\_abl1\_bootstrap\_cis\_verified.csv}, \texttt{egfr\_abl1\_permutation\_tests.csv}.}\label{fig3}
\end{figure}

\begin{figure}[!htb]
\centering
\includegraphics[width=\textwidth]{figure5_egfr_abl1_labeled_scatter.pdf}
\caption{EGFR/ABL1: mutation-labeled scatterplots of zero-shot score versus document-matched $\Delta\mathrm{p}X$, for ESM-1b and ESM-2-150M (dashed line: unweighted linear fit, shown for visual reference only, not implying a linearity assumption). At $n=6$--9, individual mutations are directly visible as high-leverage points; combined with the leave-one-out analysis in \S\ref{subsec2-3}, these plots make the sample-size fragility of every EGFR/ABL1 correlation reported in this study inspectable rather than only summarized. Source data: \texttt{egfr\_abl1\_final\_results\_document\_matched.csv}.}\label{figScatter}
\end{figure}

\FloatBarrier

\subsection{Sensitivity analyses and the MET replication gap}\label{subsec2-4}

Because ESM-1b's 1024-token architectural limit required scoring EGFR and ABL1 within truncated kinase-domain windows rather than at full length, we tested whether this truncation affected the results. ESM-2, which has no such length limit, was scored on a 17-mutation non-cascade subset both within the truncated window and at full protein length; the two scores agreed almost perfectly (Spearman $\rho = 0.998$, Pearson $r = 0.999$; Figure~\ref{fig4}). ESM-1b itself was rescored under two alternative window sizes (100- and 250-residue flanks, versus the original 150-residue flank); agreement remained strong in both cases ($\rho = 0.934$ and $\rho = 0.966$ respectively), with the modestly lower agreement relative to ESM-2 consistent with ESM-1b's use of learned absolute positional embeddings (window-boundary sensitive by construction) versus ESM-2's rotary embeddings (position-shift invariant). Window choice does not materially affect the results reported here.

We also tested whether the specific zero-shot scoring convention explained the MET replication gap (\S\ref{subsec2-1}). Wildtype-marginal scoring (a single unmasked forward pass, reading the model's probability at each position directly) achieved $\rho = 0.9305$ against Estevam et al.'s published scores, compared to $\rho = 0.9199$ for the masked-marginal method used throughout the rest of this study. The scoring-convention choice therefore explains only a small fraction of the gap to our a priori target; we retained masked-marginal scoring for all main results, both because it is the method Estevam et al.\ themselves report using and because it is generally the better-performing convention in the broader zero-shot variant-effect literature \citep{Meier2021}, independent of this specific replication question.

\begin{figure}[!htb]
\centering
\includegraphics[width=\textwidth]{figure4_sensitivity_analysis.pdf}
\caption{Window/flank sensitivity analysis on a 17-mutation non-cascade subset. Left: ESM-2-150M scored within the truncated kinase-domain window versus at full protein length ($\rho=0.998$). Center and right: ESM-1b scored under the original 150-residue flank versus 100- and 250-residue alternatives ($\rho=0.934$ and $\rho=0.966$ respectively).}\label{fig4}
\end{figure}

\FloatBarrier

\section{Methods}\label{sec3}

\subsection{Datasets}\label{subsec3-1}

\textbf{MET.} We used the deep mutational scanning dataset of Estevam et al.\ \citep{Estevam2025} (\texttt{fraser-lab/\allowbreak MET\_\allowbreak kinase\_\allowbreak Inhibitor\_\allowbreak DMS}, MIT license), comprising 5{,}764 kinase-domain variants assayed against nine ATP-competitive inhibitors (cabozantinib, capmatinib, crizotinib, glesatinib, glumetinib, merestinib, savolitinib, tepotinib, and an NVP compound), of which 5{,}434 variants had complete scores across all models after quality control. The wild-type construct sequence (287 residues) was reconstructed from the authors' deposited data and independently verified against the UniProt canonical MET sequence.

\textbf{EGFR and ABL1.} Canonical sequences were retrieved from UniProt \citep{UniProt2025} (EGFR: P00533; ABL1: P00519, isoform 1b, matching clinical mutation numbering). Kinase-domain boundaries (EGFR: residues 712--979; ABL1: residues 242--493) were retrieved via the UniProt domain-feature API. Because ESM-1b's learned positional embeddings impose a hard 1024-token limit, both full-length sequences (EGFR: 1210 aa; ABL1: 1130 aa) exceeded the usable length; we constructed scoring windows by extending each kinase-domain boundary by a 150-residue flank in both directions, clipped to sequence bounds and to always include all mutation positions (EGFR window: residues 562--1129, 568 aa; ABL1 window: residues 92--643, 552 aa), applied identically across all four scored models.

Resistance-mutation panels were assembled from the clinical and DMS literature. For EGFR, ChEMBL querying revealed that compound mutations (e.g.\ L858R+T790M+C797S) vastly outnumber isolated single mutations in the bioactivity data, reflecting that resistance mutations arise sequentially on a clinical background; we therefore constructed a background-aware cascade (WT $\rightarrow$ L858R $\rightarrow$ L858R+T790M $\rightarrow$ L858R+T790M+C797S) with each step scored against its immediate predecessor rather than against wild-type. Three additional, independently occurring EGFR activating mutations (L861Q, G719C, G719S) were scored against wild-type directly. For ABL1, which showed no compound-mutation pattern in ChEMBL, 13 clinically documented resistance mutations were scored against wild-type (M244V, G250E, Y253F, Y253H, E255K, E255V, T315I, M351T, F359V, F317L, F317I, H396P, Q252H).

\subsection{Zero-shot scoring}\label{subsec3-2}

We used masked-marginal scoring \citep{Meier2021}: for each mutation position, the wild-type (or background) sequence token at that position is masked, the model's log-probability distribution over amino acids at the masked position is computed via a forward pass, and the zero-shot score is $\log P(\text{mut} \mid \text{context}) - \log P(\text{wt} \mid \text{context})$. We scored ESM-1b (650M parameters) and ESM-2 (150M, 650M, and 3B parameter checkpoints), using identical masked-marginal scoring, identical sequence windows, and identical mutation positions across all four models for a given kinase.

\subsection{Non-learned baselines}\label{subsec3-3}

\textbf{BLOSUM62.} Substitution score for the wild-type-to-mutant amino acid pair \citep{Henikoff1992}.

\textbf{Structural distances.} Two distances were computed from real inhibitor-bound crystal structures rather than approximated: distance to a reference ligand, and distance to the validated ATP-binding pocket. Structures were retrieved from the RCSB Protein Data Bank \citep{Burley2021}. For EGFR we used PDB 6LUD (L858R/T790M/C797S in complex with osimertinib) \citep{Kashima2020}, chosen because its engineered mutations match our cascade panel exactly; for ABL1 we used PDB 1IEP (wild-type in complex with imatinib) \citep{Nagar2002}. Both structures' residue numbering was validated directly against four literature-established, non-mutated catalytic residues per kinase before use (EGFR: K745, E762, M793, D855; ABL1: K271, E286, M318, D381); for 6LUD, the three engineered mutation sites were additionally confirmed to show the expected mutant (not wild-type) identity, verifying both the numbering convention and that the correct structure had been retrieved. We deliberately name this baseline \emph{distance to reference ligand} rather than ``distance to the inhibitor,'' because each mutation's document-matched $\Delta\mathrm{p}X$ outcome (\S\ref{subsec3-4}) is aggregated across every compound with a qualifying ChEMBL record for that mutation, typically several to dozens of distinct compounds, whereas the crystallographic distance is computed to a single bound ligand (osimertinib or imatinib) in one selected structure. It is therefore a fixed reference-geometry feature, not a compound-specific distance matched to each individual measurement it is being correlated against; computing a genuinely compound-specific version would require either a separate crystal or predicted structure for every compound in the matched-pair set, or docking, neither of which we attempted here. Distance to reference ligand and distance to ATP pocket were each computed as the minimum heavy-atom distance from the mutated residue to the bound-ligand atoms or to the four validated catalytic-residue atoms, respectively -- the same distance metric applied to two different reference atom sets, keeping the two baselines directly comparable to each other, though not to a per-compound measurement.

For MET, we reused the structural and biophysical features (distance to ATP, distance to reference ligand, pocket volume, hydrophobicity, polarity, and stability estimates) provided directly by Estevam et al.\ \citep{Estevam2025} rather than recomputing them, since they had already been validated by the original authors; we did not independently verify whether their own distance-to-ligand feature was computed per treatment inhibitor or from a single reference structure, and note this as an open question rather than assuming either.

\subsection{ChEMBL matched-pair construction}\label{subsec3-4}

Bioactivity records (IC$_{50}$, K$_i$, K$_d$) were retrieved from ChEMBL \citep{Gaulton2012,Zdrazil2024} for each kinase and mutation via the \texttt{assay\_variant\_mutation} field, restricted to human-target records. Values were converted to a common $\mathrm{p}X = 9 - \log_{10}(\text{value in nM})$ scale. Wild-type and mutant records were matched by compound identity (\texttt{molecule\_chembl\_id}), requiring agreement on assay type (IC$_{50}$/K$_i$/K$_d$), source publication (\texttt{document\_chembl\_id}), and assay format (\texttt{bao\_format}) -- a stricter criterion than compound-and-assay-type alone, introduced after identifying that the looser criterion could pair records from substantially different assay protocols for the same compound (\S\ref{subsec2-3}). We confirmed no stricter same-\texttt{assay\_chembl\_id} matches existed for any mutation, consistent with ChEMBL assay entries being curated per target-variant rather than as head-to-head comparisons. Within each matching compound-document-format group, the median $\mathrm{p}X$ was taken on each side; the resulting per-compound $\Delta \mathrm{p}X$ values were then aggregated to one value per mutation by taking the median across compounds, so that no single heavily-studied compound could dominate a mutation's estimated effect size. Mutations with fewer than ten matched compounds were flagged as unreliable and excluded from correlation analyses, though retained in the full results table.

\subsection{Statistics}\label{subsec3-5}

All correlations were computed as Spearman rank correlations, given no assumption of linearity and small-to-moderate sample sizes. For MET, per-drug correlations were computed separately for each of the nine inhibitors (rather than pooling across drugs before correlating) and combined via Fisher-z-transformed averaging (each per-drug $\rho$ transformed via $\operatorname{arctanh}$, averaged, and transformed back via $\tanh$; correlations were clipped to $[-0.999, 0.999]$ before transformation to avoid unbounded values); model-versus-model comparisons used a paired Wilcoxon signed-rank test \citep{Wilcoxon1945} across the nine drugs, which we note treats the nine per-drug correlations as independent even though all nine share the same underlying 5{,}434-mutation panel (\S\ref{subsec2-2}, \S\ref{subsec4-1}).

Two distinct bootstrap procedures were used for MET, and we document both explicitly because they differ in a way that matters for how their results should be read. The shared-mutation cluster bootstrap (ESM-1b-versus-ESM-2 comparison, \S\ref{subsec2-2}) draws, on each of 2{,}000 iterations, one set of mutation indices sampled with replacement from the full 5{,}434-mutation panel; that same index set is applied to reconstruct all nine drug-specific subsets for all four models simultaneously, so ESM-1b and every ESM-2 size see identical resampled mutations within an iteration. Rows with a missing score for the model or drug in question are dropped (complete-case analysis) before each per-drug correlation is computed; a drug-model combination is excluded from that iteration's Fisher-z average if fewer than 10 mutations remain after this filtering or if either variable is constant within the resample. The nine resulting per-drug correlations are Fisher-z-averaged per model per iteration as described above, and the ESM-1b-versus-ESM-2 difference in this Fisher-z-averaged value is recorded; the reported point estimate is the mean of this difference across all 2{,}000 iterations, the 95\% confidence interval is the 2.5th/97.5th percentile of the same distribution, and the two-sided bootstrap $p$-value uses the finite-resampling (add-one) correction $p = 2 \times \min\!\left(\frac{1 + \#(\Delta \leq 0)}{B+1}, \frac{1 + \#(\Delta \geq 0)}{B+1}\right)$, capped at 1, where $B=2{,}000$ and $\#(\cdot)$ counts iterations satisfying the condition; this avoids the invalid $p=0$ that the uncorrected empirical-proportion formula can return when no bootstrap difference crosses zero, and gives a minimum attainable two-sided $p$-value of approximately $0.001$ at this $B$. Because mutations are resampled while the nine drug conditions are held fixed, this procedure quantifies uncertainty arising from the finite mutation panel, not from the broader population of possible drugs (\S\ref{subsec2-2}).

The MET baseline-comparison bootstrap (ESM-1b-versus-non-learned-baseline comparisons, \S\ref{subsec2-2}, Figure~\ref{fig2}) uses the identical shared-mutation cluster bootstrap procedure described above for the ESM-1b-versus-ESM-2 comparison, substituting each baseline predictor for the ESM-2 model: on each of 2{,}000 iterations, one shared set of mutation indices is drawn with replacement from the full panel and applied identically to ESM-1b and to the baseline under comparison across all nine drugs; both predictors' nine per-drug correlations are Fisher-z-averaged within that iteration, and the ESM-1b-minus-baseline difference in this Fisher-z-averaged value is recorded once per iteration. The reported point estimate, 95\% CI, and two-sided bootstrap $p$-value use the same formulas as above, and Benjamini-Hochberg correction \citep{BenjaminiHochberg1995} was applied across the eight baseline comparisons.

For EGFR/ABL1, 95\% confidence intervals were computed via percentile-method bootstrap (20{,}000 resamples, vectorized rank-based Spearman computation), resampling mutations with replacement within each grouping. Significance was assessed via exact permutation tests where feasible (full enumeration of all $n!$ mutation-label permutations: 720 for EGFR, $n=6$; 362{,}880 for ABL1, $n=9$) and via Monte Carlo permutation (200{,}000 permutations) for the pooled Combined grouping ($n=15$, where exact enumeration is infeasible), rather than the asymptotic Spearman $t$-approximation, which is unreliable at these sample sizes. Model-versus-model and model-versus-baseline comparisons for EGFR/ABL1 used a paired bootstrap on the difference in Spearman $\rho$ (the same resampled mutation set used for both predictors within each iteration, 20{,}000 resamples), applied to all 6 predictor pairs involving ESM-1b within each grouping (ESM-1b versus each of the other 6 predictors; 18 comparisons total across the 3 groupings), since EGFR/ABL1 mutations are not tested against a fixed, shared compound panel and therefore do not support MET's per-drug pairing structure directly. A leave-one-out (jackknife) analysis recomputed each correlation with each single mutation removed in turn, to quantify how much a single observation can change a given estimate at these sample sizes. Benjamini-Hochberg correction \citep{BenjaminiHochberg1995} was applied to the permutation-based $p$-values across the full battery of 21 EGFR/ABL1 correlation tests (7 predictors $\times$ 3 groupings); ABL1 and EGFR are treated as the primary within-kinase groupings and Combined as a secondary pooled sensitivity check (\S\ref{subsec2-3}). A post-hoc power analysis (bivariate-normal simulation, Spearman test, $\alpha=0.05$, 20{,}000 replicates per cell, sample sizes $n=6$ to $n=100$, target correlations $\rho=0.15$ to $0.50$) quantified statistical power at the available sample sizes and estimated the sample size needed for 80\% power via both simulation and the closed-form Fisher-$z$ approximation of Bonett and Wright.

\subsection{Reproducibility}\label{subsec3-6}

All sequence retrieval, mutation-panel construction, and window-boundary numbering were validated by direct assertion against independently retrieved reference data (UniProt canonical sequences; crystal-structure catalytic-residue identities) rather than assumed. Analysis code and processed result tables are available from the corresponding author upon request.

\section{Discussion}\label{sec4}

The central empirical result of this study is a contrast in what could and could not be established, not a contradiction between kinases. For MET, a deep-mutational-scanning design provided 5{,}434 variants against nine inhibitors, and within that well-powered setting the answer to our motivating question is unambiguous: greater ESM-2 scale does not improve, and does not match, ESM-1b's zero-shot correlation with drug-specific resistance, even though scale does improve performance monotonically within the ESM-2 family considered alone. This is a genuinely non-obvious finding. ESM-2 outperforms ESM-1b on general protein-fitness and structure-prediction benchmarks \citep{Lin2023}; a model that wins on those broader benchmarks can lose on this narrower, more clinically specific task. Model generation and parameter count are not interchangeable proxies for ``better'' when the target quantity is drug-specific rather than general evolutionary fitness.

For EGFR and ABL1, the same pipeline -- including the same non-learned structural baselines validated against real crystal structures, exact permutation tests rather than asymptotic approximations, leave-one-out influence checks, and explicit paired comparisons of ESM-1b against every other predictor individually -- could not resolve the question at all, for any of the seven predictors tested, once corrected for the number of comparisons made. We do not read this as evidence against the MET finding, nor as evidence that PLM scale is irrelevant to EGFR/ABL1 resistance specifically. We read it as evidence of a structural asymmetry between the two data sources available for this class of study, a straightforward consequence of how each data type is generated rather than a new empirical discovery in its own right, but one worth stating explicitly rather than leaving implicit. Deep mutational scanning provides many mutations tested against few drugs; curated bioactivity databases provide few clinically characterized mutations (by definition -- these are the mutations rare and consequential enough to have been studied and reported) tested against many drugs. A 13-to-19-mutation panel, even after literature-directed expansion and even with document- and format-matched (rather than fully controlled head-to-head) wild-type/mutant comparisons, does not carry the statistical power that a several-thousand-variant DMS panel does, regardless of how carefully the comparison is constructed -- our post-hoc power analysis (\S\ref{subsec2-3}) puts a concrete number on this: 8--13\% power at the available sample sizes to detect a MET-sized effect, against roughly 117 mutations needed for conventional 80\% power. This limits how confidently a cross-kinase generalization claim can be made from bioactivity-database-derived panels alone.

Two patterns in the EGFR/ABL1 data, though not statistically confirmed, are worth flagging as directions for adequately powered future work rather than as findings of this study. First, the one nominally significant (pre-correction) EGFR result favored ESM-2 over ESM-1b, opposite to the MET ranking; if this reversed pattern held up under greater power, it would suggest the MET finding is not universal across kinases, which would itself be a substantive and publishable result. Second, ABL1's structural distance-to-pocket baseline numerically exceeded every PLM tested, consistent with the geometric intuition that resistance mutations cluster near the binding site by definition; a well-powered demonstration that a simple, non-learned structural feature matches or beats large PLMs would be an important and more actionable finding for practical resistance-liability screening than a PLM-superiority result would be.

This study's negative result for scale should not be read as a verdict on the broader family of unsupervised or generative variant-effect approaches. Alignment-based generative models such as DeepSequence, EVmutation, and EVE \citep{Riesselman2018,Hopf2017,Frazer2021}, and retrieval-augmented autoregressive models such as Tranception \citep{Notin2022}, use architecturally distinct mechanisms from ESM's masked-language-model objective and were not tested here; whether any of these captures drug-specific resistance signal that scale-matched ESM checkpoints do not is a distinct empirical question outside this study's scope. Similarly, our structural baselines deliberately used experimentally resolved crystal structures rather than predicted ones, so this study makes no claim about whether AlphaFold-predicted structures \citep{Jumper2021} or structure-and-sequence classifiers such as AlphaMissense \citep{Cheng2023} would perform comparably; testing predicted-structure baselines against the same resistance panels, and thereby introducing structure-prediction uncertainty as an explicit variable, is a natural next step once bioactivity panels are large enough to support it.

\subsection{Limitations}\label{subsec4-1}

First, although we controlled for assay heterogeneity more strictly than a compound-and-assay-type match alone (\S\ref{subsec3-4}), even document-and-format-matched ChEMBL records are not equivalent to a single controlled head-to-head experiment: they do not guarantee an identical assay protocol, cell or biochemical system, protein construct, incubation time, or laboratory conditions between the wild-type and mutant measurements, and no stricter same-assay match existed in the data to test against. Second, the EGFR/ABL1 mutation panel, while literature-directed rather than arbitrary, remains a convenience sample of clinically documented mutations rather than a systematic saturation panel; our power analysis (\S\ref{subsec2-3}) shows that expanding it toward roughly 100 mutations is the most direct way to make this comparison adequately powered, but doing so was constrained in this study by the practical availability of matched bioactivity data. Third, the distance-to-reference-ligand baseline (\S\ref{subsec3-3}) is computed from a single crystal structure per kinase and is not compound-specific, even though the outcome it is compared against aggregates many compounds per mutation; we have renamed and scoped this baseline accordingly rather than treating it as a true per-compound structural feature. Fourth, the MET replication gap (\S\ref{subsec2-1}) is only partially explained; scoring-convention choice (\S\ref{subsec2-4}) accounts for a small fraction of it, and construct sequence, residue numbering, model checkpoint/library version, floating-point precision, sequence preprocessing, and mutation-filtering criteria remain untested candidate explanations. Fifth, BRAF was considered and excluded from this study during scope-finalization, because clinical resistance to BRAF inhibitors is driven predominantly by pathway reactivation (RAS mutation, receptor tyrosine kinase upregulation, aberrant splicing) rather than second-site kinase-domain mutations, and the rare exceptions reported in the literature had no usable ChEMBL activity data; this is a considered exclusion rather than an oversight, but it does mean this study's claims are scoped to tyrosine kinases with second-site-mutation-driven resistance mechanisms, not to kinase inhibitor resistance in general. Sixth, the MET shared-mutation cluster bootstrap (\S\ref{subsec2-2}, \S\ref{subsec3-5}) holds the panel of nine MET inhibitors fixed and resamples only mutations; it therefore quantifies uncertainty arising from the finite mutation panel for these nine specific inhibitors, not from the broader population of possible ATP-competitive MET inhibitors or of kinase inhibitors generally, and the resulting significance claims should be read as scoped to this nine-drug panel.

The public processed feature tables for MET contained 5{,}434 quality-controlled variants. The reduction from the original 5{,}764 variants occurred in the source study's upstream ROSACE quality-control pipeline \citep{Estevam2025}, which the original authors report as removing variants failing a mean-read-count threshold (mean count $\geq 4$ across replicates and timepoints), an excess-zero-count filter (more than 10 of 12 replicate/timepoint counts equal to zero), and a baseline-timepoint zero-count filter (more than 2 of 3 T0 replicate counts equal to zero). Because raw replicate-level count files were unavailable to us, we could not independently reconstruct per-filter exclusion counts; consequently, our MET analyses inherit the source study's quality-controlled variant panel rather than an independently re-audited one. Processed data and analysis code are currently available from the corresponding author upon reasonable request rather than deposited in a public, DOI-assigned repository; depositing the analysis notebooks, processed mutation panels, exact package/model-checkpoint versions, ChEMBL query date and release version, and figure source data in a citable public repository before submission is planned.

\section{Conclusion}\label{sec5}

For MET specifically, in a well-powered deep-mutational-scanning comparison, zero-shot protein language models carry real, measurable signal about drug-specific kinase-inhibitor resistance, and greater model scale within the ESM-2 family does not improve on the older, smaller ESM-1b for this specific task, even as it improves performance on general benchmarks. A shared-mutation cluster bootstrap confirms ESM-1b's advantage over every ESM-2 size and every one of eight non-learned baselines, with 95\% confidence intervals excluding zero for every comparison (Benjamini-Hochberg-adjusted $p \leq 0.001$ against every baseline; scope caveats in \S\ref{subsec4-1}), and the same untreated-versus-treated fitness gap Estevam et al.\ reported for ESM-1b alone holds, with a bootstrap confidence interval excluding zero, for all four model sizes we tested, indicating that the untreated-versus-treated contrast is shared across model families rather than being unique to ESM-1b. This MET-specific result should not be read as an established, general claim about PLMs and drug-specific resistance beyond MET: whether it generalizes across kinases remains genuinely open. Our EGFR and ABL1 analyses, built with the same rigor and the same baselines as the MET analysis, are honestly underpowered to answer that question with currently available bioactivity data, a limitation of statistical power rather than of the pipeline (\S\ref{subsec2-3}). Future work extending this comparison to additional kinases with either larger curated bioactivity panels or dedicated deep mutational scans against multiple inhibitors, while addressing the limitations identified in \S\ref{subsec4-1}, is needed before a confident cross-kinase claim -- in either direction -- can be made.

\backmatter

\bmhead{Acknowledgements}
Not applicable.

\section*{Declarations}

\begin{itemize}
\item \textbf{Funding:} Not applicable.
\item \textbf{Conflict of interest/Competing interests:} The author declares no competing interests.
\item \textbf{Ethics approval and consent to participate:} Not applicable. This study used only publicly available, previously published sequence, structural, and bioactivity data; no new experiments on human or animal subjects were performed.
\item \textbf{Consent for publication:} Not applicable.
\item \textbf{Data availability:} MET deep mutational scanning data are available from \texttt{fraser-lab/\allowbreak MET\_\allowbreak kinase\_\allowbreak Inhibitor\_\allowbreak DMS} (GitHub, MIT license). EGFR and ABL1 sequence data are available from UniProt (accessions P00533, P00519). Bioactivity data are available from ChEMBL. Crystal structures are available from the RCSB Protein Data Bank (accessions 6LUD, 1IEP). Processed result tables are available from the corresponding author upon reasonable request.
\item \textbf{Materials availability:} Not applicable.
\item \textbf{Code availability:} Analysis notebooks are available from the corresponding author upon reasonable request.
\item \textbf{Author contribution:} Mohamed Samiulla Khan conceived the study, performed all analyses, and wrote the manuscript.
\end{itemize}

\bibliography{sn-bibliography}%

\end{document}
